\section{Discussion and limitations}
\label{sec:discussion}

This section interprets the framework's claims and their boundaries. The
confirmatory analysis will determine which specific effects are observed;
the discussion is structured so that each claim is reported with its
prereqistered threshold and its predeclared fallback, rather than
reinterpreted after seeing results.

\subsection{Component attribution}

The factorial ablation (R2-A) is designed to attribute gains across three
factors: calibrated R1 belief, active verification, and process supervision.
The principal question is whether the adaptive controller improves outcomes
primarily by selecting better verification targets, by timing interventions
earlier, or by choosing more appropriate recovery actions, and whether the
training benefit comes from the process labels carrying recovery-specific
information or from the denser training gradient of step-level signals. The
critical interaction is active verification $\times$ process supervision:
if the benefit of process-supervised training is concentrated in
trajectories where the policy actively selected evidence, this supports the
hypothesis that the training signal teaches evidence acquisition rather
than only generic task behavior. If the interaction is absent, the two
components are reported as independent contributions.

\subsection{Sensitivity to cost-model choices}

The controller's decisions depend on the cost weights $\lambda_v$,
$\lambda_{\text{compute}}$, $\lambda_{\text{human}}$, and
$\lambda_{\text{damage}}$, which encode value judgments about the relative
importance of verification effort, compute, human time, and irreversible
damage. The primary result is reported at the preregistered cost
configuration, and a sensitivity analysis is run over a declared grid of
alternative weight vectors. If the adaptive advantage disappears or reverses
in a region of the cost space, that region is reported explicitly rather
than hidden by a single operating point. Users who adopt the controller
with their own cost weights are expected to re-evaluate the primary results
under those weights.

\subsection{Reward-channel independence}

The verifier-substitution analysis reports whether the policy's performance
degrades when evaluated with a different verifier implementation on the same
trajectories. The reward-channel audit reports reward--evidence disagreement:
cases where the training reward signal is positive but the independent
adjudicator rejects the outcome. A policy that maintains its performance
under verifier substitution and whose reward channel shows no leakage is
interpreted as learning scientific recovery rather than verifier-specific
artifacts. If either audit fails, the affected claim is narrowed and the
reward-hacking risk is reported prominently.

\subsection{External validity}

All results are conditional on the sampled models, task families, fault
mechanisms, and cost configurations. The controller is trained on R1
development traces and evaluated on R1 confirmatory held-out groups. The
external-validity L4 arm tests whether the controller value transfers to
naturally occurring failures, and the environment-shift evaluation set
tests whether the clean-task non-inferiority result holds under a changed
runtime or model version. Any evaluation set where the benefit disappears
is reported rather than averaged away.

\subsection{Failure cases}

The analysis reports cases where the adaptive policy made worse decisions
than a fixed baseline: over-verification on low-risk trajectories,
under-verification on rapidly propagating faults, inappropriate recovery
actions, and unnecessary abstention. The question is whether these failures
are systematic (predictable from the belief state) or stochastic. Systematic
failure modes are candidates for protocol-level fixes in a future release;
stochastic failures are characterized by their frequency and cost.

\subsection{Security and safety implications}

A policy trained to verify and recover could be used to reduce the cost of
operating unreliable agents rather than to improve reliability. The
clean-task regression gate, the verifier-substitution audit, the reward
channel audit, and the requirement for independent adjudication are
structural defenses against this risk, but they are not guarantees. A
deployed version should maintain the separation between the training reward
channel and the operational safety monitor, and should re-run the clean-task
and verifier-substitution gates under the deployment distribution.

\subsection{Boundary of claims}

H3 and H4 each require specific evidence thresholds
(Section~\ref{sec:experiments}). If either threshold is not met, the
corresponding claim is narrowed or withdrawn as specified in the
preregistered fallback. Claims are bounded to the sampled models, tasks,
fault regimes, and budget configurations; they do not extend to
substantially different agent architectures, scientific domains, or cost
structures without further evidence. The process-supervised result, if
positive, is reported as an incremental contribution over matched
outcome-only training at matched data and compute, not as a claim that
process supervision is universally superior.
